\clearpage
\chapter{MARCO APLICATIVO}

Para el presente capítulo es necesaria la lectura de los capítulos anteriores, ya que se utilizan los conceptos de sistemas distribuidos, patrón \textit{Transactional Outbox}, idempotencia y las tecnologías descritas en el marco teórico. En este capítulo se detalla el diseño, la arquitectura y la implementación de la librería propuesta.

\section{ARQUITECTURA DE LA LIBRERÍA}

La librería se diseñó siguiendo principios de modularidad, extensibilidad y bajo acoplamiento, permitiendo su integración en cualquier proyecto Spring Boot con mínima configuración.

\subsection{Estructura Modular}

La librería se compone de tres módulos principales:

\begin{itemize}
    \item \textbf{Módulo de Persistencia}: responsable de almacenar los eventos en la tabla Outbox dentro de la misma transacción de negocio.
    \item \textbf{Módulo Publicador}: encargado de leer periódicamente los eventos pendientes y publicarlos en Apache Kafka.
    \item \textbf{Módulo de Idempotencia}: permite a los consumidores detectar y descartar eventos duplicados.
\end{itemize}

La Figura~\ref{fig:arquitectura-libreria} presenta la arquitectura modular de la librería.

\begin{figure}[H]
    \centering
    \begin{tikzpicture}[
        node distance=1.15cm and 1.25cm,
        every node/.style={font=\small}
    ]
        \node[tesisService, minimum width=4.0cm] (app) {Aplicación Spring Boot};
        \node[tesisProcess, below=of app, minimum width=4.7cm] (lib) {Librería Transactional Outbox};

        \node[tesisBox, fill=tesisSoftBlue, below left=1.25cm and 2.05cm of lib, minimum width=2.6cm] (persist) {Módulo\\Persistencia};
        \node[tesisBox, fill=tesisSoftAmber, below=1.25cm of lib, minimum width=2.6cm] (publisher) {Módulo\\Publicador};
        \node[tesisBox, fill=tesisSoftGreen, below right=1.25cm and 2.05cm of lib, minimum width=2.6cm] (idempotent) {Módulo\\Idempotencia};

        \node[tesisDatabase, below=of persist, minimum width=2.6cm] (db) {PostgreSQL};
        \node[tesisBroker, below=of publisher, minimum width=2.6cm] (kafka) {Apache Kafka};
        \node[tesisDatabase, below=of idempotent, minimum width=2.6cm] (processed) {Eventos\\procesados};

        \draw[tesisArrow] (app) -- (lib);
        \draw[tesisArrow] (lib) -- (persist);
        \draw[tesisArrow] (lib) -- (publisher);
        \draw[tesisArrow] (lib) -- (idempotent);
        \draw[tesisArrow] (persist) -- (db);
        \draw[tesisArrow] (publisher) -- (kafka);
        \draw[tesisDashedArrow] (idempotent) -- (processed);
    \end{tikzpicture}

    \caption{Arquitectura modular de la librería propuesta.}
    \label{fig:arquitectura-libreria}
\end{figure}


\subsection{Modelo de Datos}

La librería utiliza dos tablas en PostgreSQL: una para almacenar los eventos pendientes de publicación y otra para registrar los eventos ya procesados por los consumidores.

\subsubsection{Tabla outbox\_events}

\begin{lstlisting}[language=SQL, caption={Script de creación de la tabla outbox\_events.}]
CREATE TABLE outbox_events (
    id              UUID PRIMARY KEY DEFAULT gen_random_uuid(),
    aggregate_type  VARCHAR(255) NOT NULL,
    aggregate_id    VARCHAR(255) NOT NULL,
    event_type      VARCHAR(255) NOT NULL,
    topic           VARCHAR(255) NOT NULL,
    payload         JSONB NOT NULL,
    status          VARCHAR(20) DEFAULT 'PENDING',
    retry_count     INTEGER DEFAULT 0,
    max_retries     INTEGER DEFAULT 5,
    created_at      TIMESTAMP DEFAULT NOW(),
    processed_at    TIMESTAMP,
    next_retry_at   TIMESTAMP DEFAULT NOW()
);

CREATE INDEX idx_outbox_status ON outbox_events(status, next_retry_at);
\end{lstlisting}

\subsubsection{Tabla processed\_events}

\begin{lstlisting}[language=SQL, caption={Script de creación de la tabla processed\_events.}]
CREATE TABLE processed_events (
    event_id        UUID PRIMARY KEY,
    processed_at    TIMESTAMP DEFAULT NOW()
);

CREATE INDEX idx_processed_at ON processed_events(processed_at);
\end{lstlisting}

\section{IMPLEMENTACIÓN DEL MÓDULO DE PERSISTENCIA}

Este módulo permite registrar eventos en la tabla Outbox dentro de la misma transacción que persiste los datos de negocio.

\subsection{Entidad OutboxEvent}

\begin{lstlisting}[language=Java, caption={Entidad JPA para la tabla Outbox.}]
@Entity
@Table(name = "outbox_events")
public class OutboxEvent {

    @Id
    @GeneratedValue(strategy = GenerationType.UUID)
    private UUID id;

    @Column(name = "aggregate_type", nullable = false)
    private String aggregateType;

    @Column(name = "aggregate_id", nullable = false)
    private String aggregateId;

    @Column(name = "event_type", nullable = false)
    private String eventType;

    @Column(nullable = false)
    private String topic;

    @Column(columnDefinition = "jsonb", nullable = false)
    private String payload;

    @Column(nullable = false)
    @Enumerated(EnumType.STRING)
    private OutboxStatus status = OutboxStatus.PENDING;

    @Column(name = "retry_count")
    private int retryCount = 0;

    @Column(name = "max_retries")
    private int maxRetries = 5;

    @Column(name = "created_at")
    private LocalDateTime createdAt = LocalDateTime.now();

    @Column(name = "processed_at")
    private LocalDateTime processedAt;

    @Column(name = "next_retry_at")
    private LocalDateTime nextRetryAt = LocalDateTime.now();
}
\end{lstlisting}

\subsection{Enumeración de Estados}

\begin{lstlisting}[language=Java, caption={Estados posibles de un evento en la tabla Outbox.}]
public enum OutboxStatus {
    PENDING,
    SENT,
    FAILED
}
\end{lstlisting}

\subsection{Repositorio de Eventos}

\begin{lstlisting}[language=Java, caption={Repositorio JPA para consultar eventos pendientes.}]
@Repository
public interface OutboxEventRepository
        extends JpaRepository<OutboxEvent, UUID> {

    @Query("SELECT e FROM OutboxEvent e " +
           "WHERE e.status = 'PENDING' " +
           "AND e.nextRetryAt <= :now " +
           "ORDER BY e.createdAt ASC")
    List<OutboxEvent> findPendingEvents(
        @Param("now") LocalDateTime now,
        Pageable pageable
    );

    @Modifying
    @Query("DELETE FROM OutboxEvent e " +
           "WHERE e.status = 'SENT' " +
           "AND e.processedAt < :before")
    int deleteProcessedBefore(
        @Param("before") LocalDateTime before
    );
}
\end{lstlisting}

\subsection{Servicio de Registro de Eventos}

\begin{lstlisting}[language=Java, caption={Servicio para registrar eventos en la tabla Outbox.}]
@Service
public class OutboxEventService {

    private final OutboxEventRepository repository;
    private final ObjectMapper objectMapper;

    public OutboxEventService(
            OutboxEventRepository repository,
            ObjectMapper objectMapper) {
        this.repository = repository;
        this.objectMapper = objectMapper;
    }

    @Transactional
    public OutboxEvent registerEvent(
            String aggregateType,
            String aggregateId,
            String eventType,
            String topic,
            Object payload) {

        OutboxEvent event = new OutboxEvent();
        event.setAggregateType(aggregateType);
        event.setAggregateId(aggregateId);
        event.setEventType(eventType);
        event.setTopic(topic);
        event.setPayload(serialize(payload));

        return repository.save(event);
    }

    private String serialize(Object payload) {
        try {
            return objectMapper.writeValueAsString(payload);
        } catch (JsonProcessingException e) {
            throw new RuntimeException(
                "Error serializando evento", e);
        }
    }
}
\end{lstlisting}

\section{IMPLEMENTACIÓN DEL MÓDULO PUBLICADOR}

% TODO: Agregar imagen - Diagrama de flujo del ciclo de vida de un evento: PENDING -> (publicación) -> SENT o PENDING -> (falla) -> reintento -> SENT/FAILED. Puede ser un diagrama de estados mostrando las transiciones.

El módulo publicador implementa la estrategia de \textit{Polling Publisher}: un proceso programado que consulta periódicamente la tabla Outbox, publica los eventos pendientes en Kafka y los marca como procesados.

\subsection{Polling Publisher}

\begin{lstlisting}[language=Java, caption={Componente de publicación periódica de eventos.}]
@Component
public class OutboxPollingPublisher {

    private final OutboxEventRepository repository;
    private final KafkaTemplate<String, String> kafkaTemplate;
    private final OutboxProperties properties;

    @Scheduled(fixedDelayString =
        "${outbox.polling.interval:1000}")
    @Transactional
    public void publishPendingEvents() {
        List<OutboxEvent> events =
            repository.findPendingEvents(
                LocalDateTime.now(),
                PageRequest.of(0,
                    properties.getBatchSize())
            );

        for (OutboxEvent event : events) {
            try {
                sendToKafka(event);
                markAsSent(event);
            } catch (Exception e) {
                handleFailure(event);
            }
        }
    }

    private void sendToKafka(OutboxEvent event) {
        kafkaTemplate.send(
            event.getTopic(),
            event.getAggregateId(),
            event.getPayload()
        ).get(5, TimeUnit.SECONDS);
    }

    private void markAsSent(OutboxEvent event) {
        event.setStatus(OutboxStatus.SENT);
        event.setProcessedAt(LocalDateTime.now());
        repository.save(event);
    }

    private void handleFailure(OutboxEvent event) {
        event.setRetryCount(
            event.getRetryCount() + 1);

        if (event.getRetryCount() >= event.getMaxRetries()) {
            event.setStatus(OutboxStatus.FAILED);
        } else {
            long delay = (long) Math.pow(2,
                event.getRetryCount());
            event.setNextRetryAt(
                LocalDateTime.now()
                    .plusSeconds(delay));
        }
        repository.save(event);
    }
}
\end{lstlisting}

\subsection{Mecanismo de Reintento con Backoff Exponencial}

Cuando la publicación de un evento falla, el sistema incrementa un contador de reintentos y calcula el próximo momento de reintento utilizando \textit{backoff exponencial}. La fórmula aplicada es:

\[
    t_{next} = t_{current} + 2^{retry\_count} \text{ segundos}
\]

De esta manera, los reintentos se espacian progresivamente: 2s, 4s, 8s, 16s, 32s. Si se supera el número máximo de reintentos configurado, el evento se marca como \texttt{FAILED} para revisión manual.

La Tabla~\ref{tab:retry_backoff} ilustra la progresión del backoff exponencial.

\begin{table}[H]
    \centering
    \caption{Progresión del backoff exponencial en reintentos.}
    \begin{tabular}{|c|c|c|}
        \hline
        \textbf{Reintento} & \textbf{Demora (s)} & \textbf{Tiempo acumulado} \\
        \hline
        1 & 2 & 2s \\
        \hline
        2 & 4 & 6s \\
        \hline
        3 & 8 & 14s \\
        \hline
        4 & 16 & 30s \\
        \hline
        5 & 32 & 62s \\
        \hline
    \end{tabular}
    \label{tab:retry_backoff}
\end{table}

\section{IMPLEMENTACIÓN DEL MÓDULO DE IDEMPOTENCIA}

Este módulo permite a los servicios consumidores detectar eventos duplicados y evitar procesarlos más de una vez.

\subsection{Entidad ProcessedEvent}

\begin{lstlisting}[language=Java, caption={Entidad JPA para registrar eventos procesados.}]
@Entity
@Table(name = "processed_events")
public class ProcessedEvent {

    @Id
    @Column(name = "event_id")
    private UUID eventId;

    @Column(name = "processed_at")
    private LocalDateTime processedAt =
        LocalDateTime.now();
}
\end{lstlisting}

\subsection{Servicio de Idempotencia}

\begin{lstlisting}[language=Java, caption={Servicio para verificar y registrar eventos procesados.}]
@Service
public class IdempotencyService {

    private final ProcessedEventRepository repository;

    public IdempotencyService(
            ProcessedEventRepository repository) {
        this.repository = repository;
    }

    @Transactional
    public boolean isAlreadyProcessed(UUID eventId) {
        return repository.existsById(eventId);
    }

    @Transactional
    public void markAsProcessed(UUID eventId) {
        ProcessedEvent record = new ProcessedEvent();
        record.setEventId(eventId);
        repository.save(record);
    }
}
\end{lstlisting}

\subsection{Anotación para Consumidores Idempotentes}

Para facilitar el uso de la idempotencia, la librería proporciona una anotación personalizada que puede aplicarse a los métodos de consumo de eventos.

\begin{lstlisting}[language=Java, caption={Anotación personalizada para consumidores idempotentes.}]
@Target(ElementType.METHOD)
@Retention(RetentionPolicy.RUNTIME)
public @interface IdempotentConsumer {
    String eventIdHeader() default "event-id";
}
\end{lstlisting}

\begin{lstlisting}[language=Java, caption={Aspecto que intercepta los consumidores anotados.}]
@Aspect
@Component
public class IdempotentConsumerAspect {

    private final IdempotencyService idempotencyService;

    @Around("@annotation(idempotentConsumer)")
    public Object checkIdempotency(
            ProceedingJoinPoint joinPoint,
            IdempotentConsumer idempotentConsumer)
            throws Throwable {

        UUID eventId = extractEventId(joinPoint,
            idempotentConsumer.eventIdHeader());

        if (idempotencyService
                .isAlreadyProcessed(eventId)) {
            return null; // Evento duplicado, descartar
        }

        Object result = joinPoint.proceed();
        idempotencyService.markAsProcessed(eventId);
        return result;
    }
}
\end{lstlisting}

\section{LIMPIEZA AUTOMÁTICA DE EVENTOS}

Para evitar el crecimiento indefinido de la tabla Outbox, la librería incluye un mecanismo de limpieza automática que elimina periódicamente los eventos ya publicados.

\begin{lstlisting}[language=Java, caption={Componente de limpieza automática de eventos procesados.}]
@Component
public class OutboxCleanupScheduler {

    private final OutboxEventRepository repository;
    private final OutboxProperties properties;

    @Scheduled(cron =
        "${outbox.cleanup.cron:0 0 2 * * ?}")
    @Transactional
    public void cleanupProcessedEvents() {
        LocalDateTime threshold =
            LocalDateTime.now().minusDays(
                properties.getRetentionDays());

        int deleted = repository
            .deleteProcessedBefore(threshold);

        log.info("Limpieza completada: {} eventos " +
            "eliminados", deleted);
    }
}
\end{lstlisting}

\section{AUTO-CONFIGURACIÓN DE SPRING BOOT}

La librería se distribuye como un \textit{Spring Boot Starter}, lo que permite su activación automática mediante la inclusión de la dependencia en el proyecto.

\subsection{Clase de Configuración}

\begin{lstlisting}[language=Java, caption={Auto-configuración de la librería.}]
@Configuration
@EnableScheduling
@EnableJpaRepositories(
    basePackages = "com.outbox.library.repository")
@EntityScan(
    basePackages = "com.outbox.library.entity")
@ConditionalOnProperty(
    name = "outbox.enabled",
    havingValue = "true",
    matchIfMissing = true)
public class OutboxAutoConfiguration {

    @Bean
    @ConditionalOnMissingBean
    public OutboxEventService outboxEventService(
            OutboxEventRepository repository,
            ObjectMapper objectMapper) {
        return new OutboxEventService(
            repository, objectMapper);
    }

    @Bean
    @ConditionalOnMissingBean
    public OutboxPollingPublisher pollingPublisher(
            OutboxEventRepository repository,
            KafkaTemplate<String, String> kafkaTemplate,
            OutboxProperties properties) {
        return new OutboxPollingPublisher(
            repository, kafkaTemplate, properties);
    }

    @Bean
    @ConditionalOnMissingBean
    public IdempotencyService idempotencyService(
            ProcessedEventRepository repository) {
        return new IdempotencyService(repository);
    }
}
\end{lstlisting}

\subsection{Propiedades de Configuración}

\begin{lstlisting}[language=Java, caption={Propiedades configurables de la librería.}]
@ConfigurationProperties(prefix = "outbox")
public class OutboxProperties {
    private boolean enabled = true;
    private int batchSize = 50;
    private long pollingInterval = 1000;
    private int maxRetries = 5;
    private int retentionDays = 7;
    private String cleanupCron = "0 0 2 * * ?";
}
\end{lstlisting}

\subsection{Configuración en application.yml}

El usuario de la librería puede personalizar el comportamiento mediante el archivo de configuración de Spring Boot:

\begin{lstlisting}[caption={Ejemplo de configuración en application.yml.}]
outbox:
  enabled: true
  batch-size: 50
  polling-interval: 1000
  max-retries: 5
  retention-days: 7
  cleanup-cron: "0 0 2 * * ?"

spring:
  kafka:
    bootstrap-servers: localhost:9092
  datasource:
    url: jdbc:postgresql://localhost:5432/mydb
\end{lstlisting}

\section{EJEMPLO DE USO}

% TODO: Agregar imagen - Diagrama de secuencia mostrando la interacción entre el Servicio Productor, la Base de Datos (tabla Outbox), el Polling Publisher, Kafka y el Servicio Consumidor (con idempotencia).

A continuación se muestra un ejemplo completo de cómo un servicio de pedidos utiliza la librería para publicar eventos de forma confiable.

\subsection{Servicio Productor}

\begin{lstlisting}[language=Java, caption={Ejemplo de uso de la librería en un servicio de pedidos.}]
@Service
public class OrderService {

    private final OrderRepository orderRepository;
    private final OutboxEventService outboxService;

    @Transactional
    public Order createOrder(CreateOrderRequest request) {
        // 1. Persistir la orden
        Order order = new Order();
        order.setCustomerId(request.getCustomerId());
        order.setTotal(request.getTotal());
        order.setStatus(OrderStatus.CREATED);
        order = orderRepository.save(order);

        // 2. Registrar evento en la tabla Outbox
        //    (misma transaccion)
        OrderCreatedEvent event =
            new OrderCreatedEvent(
                order.getId(),
                order.getCustomerId(),
                order.getTotal());

        outboxService.registerEvent(
            "Order",
            order.getId().toString(),
            "OrderCreated",
            "orders.events",
            event
        );

        return order;
    }
}
\end{lstlisting}

\subsection{Servicio Consumidor con Idempotencia}

\begin{lstlisting}[language=Java, caption={Ejemplo de consumidor idempotente.}]
@Service
public class InventoryConsumer {

    private final InventoryService inventoryService;

    @KafkaListener(topics = "orders.events")
    @IdempotentConsumer(eventIdHeader = "event-id")
    public void handleOrderCreated(
            ConsumerRecord<String, String> record) {

        OrderCreatedEvent event =
            deserialize(record.value());

        inventoryService.reserveStock(
            event.getOrderId(),
            event.getItems()
        );
    }
}
\end{lstlisting}

\section{ESTRUCTURA DEL PROYECTO}

La Tabla~\ref{tab:estructura_proyecto} presenta la estructura de paquetes de la librería.

\begin{table}[H]
    \centering
    \caption{Estructura de paquetes de la librería.}
    \begin{tabular}{|l|p{7cm}|}
        \hline
        \textbf{Paquete} & \textbf{Responsabilidad} \\
        \hline
        \texttt{entity} & Entidades JPA (OutboxEvent, ProcessedEvent). \\
        \hline
        \texttt{repository} & Repositorios Spring Data JPA. \\
        \hline
        \texttt{service} & Servicios de negocio (registro, idempotencia). \\
        \hline
        \texttt{publisher} & Polling Publisher y lógica de reintentos. \\
        \hline
        \texttt{cleanup} & Scheduler de limpieza automática. \\
        \hline
        \texttt{config} & Auto-configuración y propiedades. \\
        \hline
        \texttt{annotation} & Anotaciones personalizadas. \\
        \hline
        \texttt{aspect} & Aspectos AOP para idempotencia. \\
        \hline
    \end{tabular}
    \label{tab:estructura_proyecto}
\end{table}

\section{DEPENDENCIAS DEL PROYECTO}

\begin{lstlisting}[caption={Dependencias principales en build.gradle.}]
dependencies {
    implementation 'org.springframework.boot:spring-boot-starter-data-jpa'
    implementation 'org.springframework.kafka:spring-kafka'
    implementation 'org.springframework.boot:spring-boot-starter-aop'
    implementation 'com.fasterxml.jackson.core:jackson-databind'

    runtimeOnly 'org.postgresql:postgresql'

    testImplementation 'org.springframework.boot:spring-boot-starter-test'
    testImplementation 'org.testcontainers:postgresql'
    testImplementation 'org.testcontainers:kafka'
}
\end{lstlisting}

